\def\stat{l1-5}





\def\tit{МЕТОД ПОВЫШЕНИЯ ОТКАЗОУСТОЙЧИВОСТИ СИСТЕМ 


ПОДДЕРЖАНИЯ КОГЕРЕНТНОСТИ КЭШ}





\def\titkol{Метод повышения отказоустойчивости систем 


поддержания когерентности кэш} 





\def\autkol{Б.\,З.~Шмейлин}





\def\aut{Б.\,З.~Шмейлин$^1$}





\titel{\tit}{\aut}{\autkol}{\titkol}





%{\renewcommand{\thefootnote}{\fnsymbol{footnote}}\footnotetext[1]


%{Работа выполнена при финансовой поддержке РФФИ


%(проект №10--07--00021) и программы ОНИТ РАН <<Информационные


%технологии и анализ сложных сис\-тем>> (проект 1.5).}}





\renewcommand{\thefootnote}{\arabic{footnote}}


\footnotetext[1]{Институт проблем информатики Российской академии наук, shmeilin@mail.ru}








\Abst{Представлен метод повышения отказоустойчивости сис\-темы поддержания 


когерентности кэш-памяти с протоколом наблюдения. В~предла\-га\-емом методе 


кодируются запросы, состояния и сигналы в схеме реализации протокола. Для достижения 


устойчивости к кратковременным отказам в логике поддержания когерентности к 


закодированным строкам запросов, состояний и сигналов применяется код, 


корректирующий одиночные и обнаруживающий двойные ошибки.}





\KW{когерентность кэш; протоколы наблюдения; отказоустойчивость; 


корректирующие коды} 








  \vskip 14pt plus 9pt minus 6pt








      \thispagestyle{headings}





      


            \label{st\stat}





\section{Введение}





      В современных микропроцессорных сис\-темах все возрастающее 


значение приобретает их отказоустойчивость. Отказоустойчивость~--- это 


свойство сис\-те\-мы, которое обеспечивает ей возможность продолжения 


действий, заданных программой, после возникновения неисправностей. 


Введение от\-ка\-зо\-устой\-чи\-вости требует избыточного аппаратного и 


программного обеспечения. В~настоящее время микропроцессорные 


сис\-темы  подвержены кратковременным отказам из-за уменьшающихся 


размеров кристаллов, плотности размещения элементов и снижения 


напряжения питания.  Кратковременные отказы (сбои), вызванные как 


различными воздействиями окружающей среды, так и внутренними 


перекрестными связями, становятся более частыми явлениями, чем отказы 


аппаратуры, обусловленные их производством.


      


      В последнее время все большее развитие получают 


многопроцессорные сис\-те\-мы (МС) как с разделенной, так и с общей 


памятью. В~этих сис\-темах каждый процессор непосредственно связан со 


своей частной кэш-памятью. При этом возникает проблема поддержания 


согласованности или когерентности кэш-памяти.  Когерентность кэш-па\-мя\-тей 


обеспечивает возможность  какому-либо процессору получить 


достоверную копию данных после модификации их другим процессором. 


Для поддержания когерентности кэш-памятей служит специальный 


механизм~--- протокол когерентности. Одним из распространенных 


протоколов является протокол наблюдения (snooping protocol). Работа по 


этому протоколу заключается в наблюдении контроллером частной 


кэш-памяти за событиями на шине общей памяти. Сбой в сис\-те\-ме поддержания 


когерентности приводит к серьезным последствиям. Так, например, наличие 


недостоверной копии строки может привести к распространению ошибки 


сразу в  нескольких процессорах.





      \begin{figure}[b] %fig1


      \vspace*{1pt}


\begin{center}


\mbox{%


\epsfxsize=113.197mm


\epsfbox{ipi-21(1)-5-1.eps}


}


\end{center}


\vspace*{-9pt}


      \Caption{Структура многопроцессорной сис\-темы


      \label{f5-1}}


      \end{figure}


      


      Данная статья посвящена разработке схемы контроллера кэш, 


реализующего алгоритм поддержания когерентности по протоколу 


наблюдения и устойчивого к ошибкам, вызванных кратковременными 


отказами. Вопросам контроля и повышения отказоустойчивости схем 


поддержания когерентности в последние годы уделяется большое 


внимание~[1--4]. В~работах~[2, 4] исследуются отказы и сбои сис\-темы с 


протоколом когерентности на основе справочника. В~[4] предлагаются 


схемы контроля и коррекции строк справочника, работа~[2] посвящена 


вопросам отказоустойчивости межпроцессорных связей в протоколах со 


справочником. В~статье~[3] рассматривается отказоустойчивая сис\-тема с 


маркерами. В~статье~[1] предлагается устройство контроля работы сис\-темы 


с протоколом наблюдения. Это устройство контроля позволяет обнаружить 


ошибки в логике поддержания когерентности, однако оно не предназначено 


для повышения отказоустойчивости схемы, реализующей эту логику. 


      


\section{Протоколы наблюдения поддержания 


когерентности кэш-памятей}





Рассмотрим структуру многопроцессорной сис\-темы, представленную на 


рис.~\ref{f5-1}. Процессор посылает запрос на чтение или запись на свою 


кэш-память. 








      При запросе на чтение и наличии валидной строки в кэш происходит 


чтение данных. Если таковая строка отсутствует в кэш-памяти, запрос через 


контроллер кэш-памяти передается на шину. Затем данные передаются из 


кэш-памяти, обладающей ими. Если такая кэш-память не находится, то 


данные передаются из главной памяти. При запросе процессора на запись в 


свою кэш-память управление  должно выдать разрешение на запись, т.\,е.\ 


разрешить конфликт между несколькими процессорами, выставившими 


запросы на запись. После получения разрешения процессор пишет в 


валидную строку, имеющуюся на данный момент в его кэш-памяти, и 


помечает эту строку как модифицированную. Если таковой  строки не 


находится, то запрос передается на шину главной памяти, так же как и при 


чтении. При получении запрашиваемой строки процессор записывает в нее 


информацию и при этом посылается запрос по шине на инвалидацию этой 


строки во всех кэш-памятях других процессоров и в главной памяти.


      


      Имеется несколько протоколов наблюдения. Мы рассмотрим в этой 


работе два наиболее распространенных протокола~--- MESI и MОESI. 


В~названиях этих протоколов использована аббревиатура из названий 


состояний. В~протоколах поддержания когерентности кэш-памяти каждой 


строке присваивается одно из рассматриваемых ниже состояний:


      \begin{itemize}


\item M~--- модифицированное, достоверная копия строки имеется только 


в данной кэш-памяти. Кэш должен затем записать строку в главную 


память перед тем, как другой процессор даст запрос на чтение. После 


записи в главную память состояние строки меняется на  Е;


\item Е~--- эксклюзивное, достоверная копия строки имеется только в 


данной кэш-па\-мя\-ти и главной памяти. При удовлетворении запроса на 


чтение другим процессором состояние строки меняется на~S;


\item  S~--- разделяемое состояние, достоверная копия строки имеется в 


главной памяти и в кэш-памятях других процессоров. При 


удовлетворении запроса на запись состояние строки меняется на~I;


\item  I~--- невалидное (недостоверное) состояние;


\item О~--- состояние владельца похоже на состояние~S, отличие его в 


том, что только строка в одной кэш-памяти может быть в состоянии~О, а 


остальные~--- в состоянии~S.


\end{itemize}





      Использование протокола MОESI  позволяет оптимизировать процесс 


поиска кэш-памяти с запрашиваемой строкой, обратившись к кэш-памяти, 


строка в которой имеет состояние~О. 


      


      При работе по этим протоколам имеются следующие запросы от 


процессора и сетевые сигналы:


      \begin{itemize}


      \item  PrRd~--- запрос на чтение от локального процессора;


\item PrWr~--- запрос на запись от локального процессора;


\item BusRd~--- сигнал и запрос на чтение от сети;


\item BusRdX~--- эксклюзивный сигнал и запрос на чтение от сети;


\item BusWB~--- сигнал обратной запись из кэш-памяти;


\item  Flush~--- сигнал и запрос на инвалидацию.


\end{itemize}





\begin{table}[b]\small


\vspace*{-6pt}


\begin{center}


\parbox{292pt}{\Caption{Запросы, переходы из одного состояния в другое и сигналы для 


протокола MESI


\label{t5-1}}





}





\vspace*{2ex}





\begin{tabular}{|c|c|c|c|}


\hline


Запрос&Исходное состояние &Конечное состояние&Сигнал\\


\hline


PrRd&M&M&---\\


PrRd&E&E&---\\


PrRd&S&S&---\\


PrRd


&I&E/S&BusRd\\


PrWr&M&M&---\\


PrWr&E&M&---\\


PrWr&S&M&Flush\\


PrWr


&I&M&BusRdX\\


BusRd&M&S&BusWB\\


BusRd&E&S&BusWB\\


BusRd&S&S&BusWB\\


BusRd


&I&I&---\\


BusRdX&M&I&BusWB\\


BusRdX&E&I&BusWB\\


BusRdX&S&I&BusWB\\


BusRdX


&I&I&---\\


Flush&M&---&SigErr\\


Flush&E&---&SigErr\\


Flush&S&I&---\\


Flush&I&I&---\\


\hline


\end{tabular}


\end{center}


\end{table}








\subsection{Протокол MESI} %2.1





      В табл.~\ref{t5-1} представлены запросы от процессора и сетевые 


запросы, переходы из одного состояния в другое и сигналы для протокола 


MESI.





Локальный запрос процессора на чтение (PrRd) в любом состоянии, кроме 


невалидного, обслуживается внутри узла и не вызывает действий сети. Такие 


запросы не изменяют состояние строки. Чтение невалидной строки (или 


строки, отсутствующей в данной кэш-памяти) приводит к запросу сети BusRd 


на чтение из памяти. Рассматриваемая здесь архитектура использует подход, 


при котором другие кэш-памяти по возможности отвечают на запрос по 


чтению. Если другие кэш-памяти не содержат запрашиваемых данных, то 


отвечает главная память. Если на запрос отвечает главная память, то строка 


кэш не представлена в других кэш, так что ее состояние становится~Е. 


В~противном случае другие кэш-памяти разделяют данные, так что 


состояние становится~S.





Локальная запись процессора (PrWr) в строку с состоянием М не меняет 


состояния и не вызывает сообщений на шине. Этот же запрос к строке с 


состоянием~Е  меняет состояние строки на~М, но не приводит к трафику по 


сети. 








      Запрос на запись PrWr в строку с состоянием S меняет состояние 


на~М и распространяет запрос по шине на инвалидацию~--- Flush, что делает 


невалидными  соответствующие строки в других кэш-памятях.


      


      Если возникает промах при локальной записи либо подлежащая 


записи строка невалидна (состояние~I), то на шину выдается сигнал 


BusRdX~--- эксклюзивный запрос на чтение от сети и инвалидацию. При 


получении валидной копии строки осуществляется запись и строка 


переходит в состояние~M.


      


      Когда контроллер кэш <<видит>> на шине запрос на чтение (BusRd) 


из име\-ющей\-ся в данной кэш-памяти валидной строки, он меняет состояние 


этой строки на~S (или оставляет строку в этом состоянии) и выдает сигнал на 


шину об обратной записи данных в главную память~--- BusWB. Если строка 


невалидна, то сигнал на шину не выдается. 


      


      Когда контроллер кэш <<видит>> на шине эксклюзивный сигнал и 


запрос на чтение (BusRdX), он меняет состояние этой строки на~I (или 


оставляет строку в этом состоянии) и  в случае валидной строки выдает 


сигнал на шину об обратной записи данных в главную память~--- BusWB. 


Если строка невалидна, то сигнал на шину не выдается. 


      


      При запросе {Flush} из сети строка переходит в состояние~I, 


если до этого она была в состоянии S, или остается в состоянии~I. При этом 


на шину сигнал не выдается. Если строка была в состоянии~М или~Е, то 


возникает сигнал об ошибке (SigErr), так как это недопустимая ситуация~--- 


запрос {Flush} посылается только для строк с состоянием S. 


      


\subsection{Протокол MOESI} %2.2





      В табл.~\ref{t5-2} представлены запросы от процессора и сетевые 


запросы, переходы из одного состояния и сигналы для протокола MESI.











      Отличия протокола MOESI от протокола MESI состоят в следующем:


      \begin{itemize}


\item при запросе на чтение процессора (PrRd) строка остается в 


состоянии~O;


\item при запросе на запись процессора (PrWr) строка из состояния~O 


переходит в состоянии~М и выдает сигнал на инвалидацию (Flush);


\item при сетевом запросе на чтение (BusRd) строка остается в 


состоянии~O;


\item при эксклюзивном запросе и запросе на инвалидацию переход из 


состояния строки~О аналогичен переходу из состояния~S. 


\end{itemize}





      Перечисленные отличия обусловлены тем, что состояние~О подобно 


состоянию~S. Исключение составляет сетевой запрос на чтение~--- BusRd, 


при котором состояние~М меняется на~О, так как запрашиваемая строка при 


таком запросе будет находиться в нескольких кэш-па\-мя\-тях и данный 


процессор  становится владельцем этой строки.





\begin{table}\small


\begin{center}


\parbox{292pt}{\Caption{Запросы, переходы из одного состояния в другое и сигналы для 


протокола MOESI


\label{t5-2}}





}





\vspace*{2ex}





\begin{tabular}{|c|c|c|c|}


\hline


Запрос&Исходное состояние &Конечное состояние&Сигнал\\


\hline


PrRd&M&M&---\\


PrRd&O&O&---\\


PrRd&E&E&---\\


PrRd&S&S&---\\


PrRd


&I&E/S&BusRd\\


PrWr&M&M&---\\


PrWr&O&M&Flush\\


PrWr&E&M&---\\


PrWr&S&M&Flush\\


PrWr


&I&M&BusRdX\\


BusRd&M&O&BusWB\\


BusRd&O&O&BusWB\\


BusRd&E&S&BusWB\\


BusRd&S&S&BusWB\\


BusRd


&I&I&---\\


BusRdX&M&I&BusWB\\


BusRdX&O&I&BusWB\\


BusRdX&E&I&BusWB\\


BusRdX&S&I&BusWB\\


BusRdX


&I&I&---\\


Flush&M&---&SigErr\\


Flush&O&I&---\\


Flush&E&---&SigErr\\


Flush&S&I&---\\


Flush&I&I&---\\


\hline


\end{tabular}


\end{center}


\end{table}


      


       \subsection{Классификация ошибок в логике поддержания 


когерентности} %2.3


      


       Следует отметить, что не все ошибки в логике реализации протокола 


когерентности приводят к искажению данных в кэш-памяти процессора, а 


лишь к потере производительности. Разделим все ошибки на три категории. 


К~первой отнесем фатальные ошибки, т.\,е.\ ошибки, которые могут 


привести к неверному выполнению программы. Ко второй категории отнесем 


ошибки, которые приводят лишь к потере производительности. И~к 


третьей~--- остальные ошибки, т.\,е.\ ошибки, не влияющие на работу 


сис\-темы.


      


      Так, к первой категории относятся такие ошибки, при которых 


состояние строки кэш-памяти вместо~I меняется на любое другое. При этом 


процессор читает недостоверную копию данных или при записи данных 


процессором не будет выдан сигнал на эксклюзивный запрос, т.\,е.\ копии 


строки в кэш-па\-мя\-тях других процессоров не будут инвалидированы. К~этой 


же группе относятся также ошибки, при которых любое состояние меняется 


на~М. При этом процессору разрешается писать в данный блок, так как он 


считает, что у него имеется единственная копия данных.


      


      Ко второй категории относятся такие, при которых состояние строки 


кэш-па\-мя\-ти вместо M, O, E или~S меняется на~I. При этом контроллер кэш, 


считая строку невалидной, запрашивает данные из кэш-памяти другого 


процессора или из главной памяти, что и приводит к потере 


производительности. К~этой же группе относятся также ошибки, при 


которых состояние~М меняется на любое другое. При этом вырабатывается 


сигнал эксклюзивный либо на инвалидацию, что приводит к 


дополнительному трафику по сети. 


      


      К третьей группе ошибок относятся такие, которые приводят, 


например, к изменению состояния~S на~E или наоборот.


      


      В табл.~\ref{t5-3} приведена полная классификация ошибок по 


категориям. В~столбцах таблицы показаны исходные состояния, в рядах~--- 


искаженные, а на пересечении столбца и колонки~--- соответствующая 


категория ошибки.


 


      \begin{table}[h]\small %3


      \vspace*{-9pt}


      \begin{center}


\parbox{196pt}{\Caption{Классификация ошибок в протоколах поддержания 


когерентности


      \label{t5-3}}


      


      }


      


      \vspace*{2ex}


      


\tabcolsep=9pt


\begin{tabular}{|c|c|c|c|c|c|}


      \hline


Исходное&\multicolumn{5}{c|}{ Ошибочное состояние}\\


\cline{2-6}


 состояние &M&O&E&S&I\\


\hline


M&X&2&2&2&2\\


O&1&X&2&2&2\\


E&1&3&X&3&2\\


S&1&3&3&X&2\\


I&1&1&1&1&X\\


\hline


\end{tabular}


\end{center}


\vspace*{-9pt}


\end{table}





\section{Отказоустойчивая схема поддержания 


когерентности кэш-памяти}





      С целью  создания отказоустойчивой схемы поддержания 


когерентности кэш-па\-мя\-ти, реализующей рассмотренные выше протоколы, 


закодируем все запросы, состояния и сигналы, а затем применим код, 


корректирующий одиночные и обнаруживающий двойные ошибки. Эти 


коды, называемые кодами SECDED (single error correction double error 


detection), широко применяются при проектировании модулей памяти~[5].


      


\subsection{Кодирование запросов, состояний и сигналов} %3.1





      \textbf{Протокол MESI.} В табл.~\ref{t5-4}--\ref{t5-6}  для протокола 


MESI приведена кодировка запросов состояний и сигналов соответственно. 





      \begin{table}[b]\small %tabl4


      \vspace*{-12pt}


      \begin{center}


      \Caption{Кодирование запросов для протокола MESI


      \label{t5-4}}


      \vspace*{2ex}


      


\tabcolsep=10pt


      \begin{tabular}{|c|c|c|}


      \hline


Запрос &Байт&Кодировка (шестнадцатеричная)\\


\hline


PrRd&0&80\\


PrWr&0&40\\


BusRd&0&20\\


BusRdX&0&10\\


Flush&0&08\\


\hline


\end{tabular}


\end{center}


%\end{table}


%      \begin{table}\small %tabl5


      \begin{center}


      \Caption{Кодирование состояний для протокола MESI


      \label{t5-5}}


      \vspace*{2ex}


      


      \begin{tabular}{|c|c|c|}


      \hline


Состояние&Байт&Кодировка (шестнадцатеричная)\\


\hline


Предыдущее &&\\


\hline


M&1&80\\


E&1&40\\


S&1&20\\


I&1&10\\


\hline


Последующее&&\\


\hline


M&1&08\\


E&1&04\\


S&1&02\\


I&1&01\\


\hline


\end{tabular}


\end{center}


%\end{table}


%      \begin{table}\small %tabl6


      \begin{center}


      \Caption{Кодирование сигналов для протокола MESI


      \label{t5-6}}


      \vspace*{2ex}


      


\tabcolsep=9.5pt


      \begin{tabular}{|c|c|c|}


      \hline


Сигнал&Байт&Кодировка (шестнадцатеричная)\\


\hline


BusRd&2&80\\


BusRdX&2&40\\


Flush&2&20\\


BusWB&2&10\\


SigErr&2&08\\


\hline


\end{tabular}


\end{center}


\end{table}








      Тогда каждая строка табл.~\ref{t5-1} может быть представлена в виде 


векторов запросов, состояний и сигналов. Эти векторы для протокола MESI 


представлены в табл.~\ref{t5-7}. Для коррекции и обнаружения ошибок в 


логике поддержания когерентности применим к векторам запросов, 


состояний и сигналов один из SECDED кодов, например код Хэмминга~[6]. 


Так как каждый вектор включает 24~бита, то применим код 31,26, т.\,е.\ код с 


26-ю информационными разрядами и 5-ю контрольными разрядами.





В табл.~7 приведены векторы запросов, состояний и сигналов для протокола MESI.


      


      Контрольные разряды формируются следующим образом. 


      


      Первый контрольный разряд~--- дополняет до четности 


информационные разряды с нечетными номерами~---  1, 3, 5 и~т.\,д. Второй 


контрольный разряд~--- дополняет до четности информационные разряды с 


номерами, у которых имеется~1 во втором разряде~--- 2, 3, 6, 7, 10 и~т.\,д. 


Третий  контрольный разряд~--- дополняет до четности информационные 


разряды с номерами, у которых имеется~1 в третьем разряде~--- 4, 5, 6, 7, 12, 


13 и~т.\,д.


      


      \begin{table}\small %tabl7


      \begin{center}


\parbox{236pt}


      {\Caption{Векторы запросов, состояний и сигналов для протокола 


MESI


      \label{t5-7}}


      


      }


      


            \vspace*{2ex}


      


      \tabcolsep=8pt


      \begin{tabular}{|c|c||c|c|}


      \hline


\multicolumn{1}{|c|}{Вектор}&Содержимое &\multicolumn{1}{c|}{Вектор}&Содержимое\\


\hline


1&808800&11&202210\\


2&804400&12&201100\\


3&802200&13&108110\\


4&801480&14&104110\\


5&408800&15&102110\\


6&404800&16&101100\\


7&402820&17&088008\\


8&401840&18&084008\\


9&208210&19&082100\\


10\hphantom{9}&204210&20&081100\\


\hline


\end{tabular}


\end{center}


\end{table}





      \begin{table}[b]\small %tabl8


      \vspace*{-12pt}


      \begin{center}


      \Caption{Кодирование состояний для протокола MОESI


      \label{t5-8}}


      \vspace*{2ex}


      


      \begin{tabular}{|c|c|c|}


      \hline


Состояние&Байт&Кодировка (шестнадцатеричная)\\


\hline


Предыдущее &&\\


\hline


M&1&80\\


E&1&40\\


S&1&20\\


I&1&10\\


О&0&04\\


\hline


Последующее&&\\


\hline


M&1&08\\


E&1&04\\


S&1&02\\


I&1&01\\


О&0&02\\


\hline


\end{tabular}


\end{center}


\end{table}





      Четвертый контрольный разряд~--- дополняет до четности 


информационные разряды с номерами, у которых имеется~1 в четвертом 


разряде~--- с~8 по 15 и~24. Пятый контрольный разряд~--- дополняет до 


четности информационные разряды с номерами, у которых имеется~1 в 


пятом разряде~--- с~16 по~24.


      


      Нарушение четности в каких-либо группах разрядов позволяет 


однозначно установить и скорректировать одиночную ошибку в любом 


информационном разряде. Если зафиксирована ошибка в незадействованных 


разрядах кода~--- разряды 6--8 байта~0, разряды 6--8 байта~2, то такая 


ошибка игнорируется. Для снижения избыточности в этих разрядах 


размещаются контрольные разряды.


      


      \textbf{Протокол MOESI.} В табл.~\ref{t5-8}  приведена кодировка 


состояний для протокола MОESI.











Векторы запросов, состояний и сигналов для протокола MОESI 


представлены в табл.~\ref{t5-9}.





\begin{table}[h]\small %tabl9


\begin{center}


\parbox{238pt}{\Caption{Векторы запросов, состояний и сигналов для протокола MОESI


\label{t5-9}}





}





\vspace*{2ex}





\tabcolsep=8pt


\begin{tabular}{|c|c||c|c|}


\hline


\multicolumn{1}{|c|}{Вектор}&Содержимое &Вектор&Содержимое \\


\hline


1&808800&14&202210\\


2&860000&15&201100\\


3&804400&16&108010\\


4&802200&17&140110\\


5&801480&18&104010\\


6&408800&19&102110\\


7&440810&20&101100\\


8&404800&21&088008\\


9&402820&22&0C0100\\


10\hphantom{9}&401840&23&084008\\


11&208210&24&082100\\


12&260010&25&081100\\


13&204210&&\\


\hline


\end{tabular}


\end{center}


\vspace*{-12pt}


\end{table}





\subsection{Реализация схемы коррекции в контроллере кэш} 


%3.2





\begin{figure} %fig2


      \vspace*{1pt}


\begin{center}


\mbox{%


\epsfxsize=115.888mm


\epsfbox{ipi-21(1)-5-2.eps}


}


\end{center}


\vspace*{-9pt}


      \Caption{Алгоритм работы контроллера кэш


\label{f5-2}}


\end{figure}





       На рис.~\ref{f5-2} представлен алгоритм работы контроллера кэш в 


части коррекции одиночных и обнаружения двойных ошибок в логике 


реализации протокола поддержания когерентности кэш. После получения 


информации о запросах, состояниях и сигналах эта информация кодируется в 


виде описанного ранее вектора. Этот вектор сравнивается с одним из 


эталонных векторов (табл.~\ref{t5-6} для протокола MESI или 


      табл.~\ref{t5-8} для протокола MОESI). При несовпадении 


формируется SECDED код. При обнаружении двойной ошибки посылается 


соответствующее сообщение. При выдаче некоторых сигналов на шину 


памяти контроллер должен получить ответ. При отсутствии ответа 


происходит <<зависание>> сис\-темы. Чтобы избежать такой ситуации, 


включается таймер и по истечении установленного времени ожидания 


выводится сообщение об отсутствии ответа. 


      


      Следует отметить, что предлагаемая схема коррекции предполагает 


определенные затраты на реализацию кода SECDED и возможно небольшую 


потерю производительности. Однако ошибки в логике когерентности могут 


привести к таким серьезным последствиям, что на обнаружение и 


корректировку потребуются большие затраты.


\section{Выводы}





\noindent


\begin{enumerate}[1.]


\item На основании анализа работы по протоколу поддержания 


когерентности кэш выявлены возможные ошибки, вызванные 


кратковременными отказами. 


\item Предложена классификация ошибок в логике реализации протокола 


когерентности кэш, согласно которой такие ошибки разделены на фатальные, 


т.\,е.\ влияющие на корректность выполнения программ, и влияющие только 


на производительность МС.


\item Для двух наиболее распространенных протоколов наблюдения~--- MESI 


и MOESI~--- предложена схема коррекции одиночных и обнаружения 


двойных ошибок с использованием кода SECDED.


\item Предложенная схема коррекции хотя и связана с введением 


избыточности в контроллере кэш и некоторой потерей производительности 


МС, однако ошибки в логике поддержания когерентности приводят к весьма 


серьезным последствиям.


\end{enumerate}





{\small\frenchspacing


{%\baselineskip=10.8pt


\addcontentsline{toc}{section}{Литература}


\begin{thebibliography}{9}





\bibitem{1-5}


\Au{Borodin D., Juurlink B.}


A low-cost cache coherence verification method for snooping 


systems~/\!\!/ 11th EUROMICRO Conference on Digital System Design 


Architectures, Methods and Tools  Proceedings, 2008. P.~219--227.








\bibitem{3-5} %2


\Au{Fernandez-Pascual R., Garcia J.\,M., Acacio~M., Duato J.}


A fault-tolerant directory-based 


cache coherence protocol for CMP architectures~/\!\!/ 


IEEE Conference (International) on Dependable Systems and Networks With FTCS and DCC Proceedings, 2008. P.~267--276.





\bibitem{4-5} %3


\Au{Fernandez-Pascual R.,  Garcia J.\,M.,  Acacio~M.\,E.,  Duato~ J.}


Extending the TokenCMP 


cache coherence protocol for low overhead fault tolerance in CMP architectures~/\!\!/ IEEE 


Trans. on Parallel and Distributed Systems Issue Date, 2008. Vol.~19. Issue~8. 


P.~1044--1056.





\bibitem{2-5} %4


\Au{Lee H.,  Cho S.,  Childers B.}


PERFECTORY: A fault-tolerant directory memory architecture~/\!\!/ 


IEEE Trans. on Computer, 2010. Vol.~59. No.\,5. P.~638--650.





\bibitem{5-5}


\Au{Ando H.}


Accelerated testing of 90~mm SPARC64V microprocessor for neutron SER~/\!\!/ 


 3rd Workshop System Effects on Logic Soft Errors Proceedings, 2007. P.~187--196.


 


 \label{end\stat}





\bibitem{6-5}


\Au{Питерсон У., Уэлдон Э.}


Коды, исправляющие ошибки.~--- М.: Мир, 1976. С.~593.


 \end{thebibliography}


}


}


